% {{{ Basic stuff

%\author{Itziar Morales Rodríguez}
%\author{Author}

% Some basic packages
%\usepackage[utf8]{inputenc}
\usepackage[T1]{fontenc}
\usepackage[main=english]{babel}
\usepackage{graphicx}
\usepackage{float}
\usepackage{booktabs}
%\usepackage{lastpage}
%\usepackage{dirtytalk}
\usepackage{float}
%\usepackage{gensymb}

% Code
%\usepackage[cachedir=\detokenize{~/.cache/minted}]{minted}
\usepackage{listings}

% Use the Roboto font
%\usepackage[sfdefault,light]{roboto}
\renewcommand{\familydefault}{\sfdefault}
\usepackage{helvet}
% Set enumerate
%\usepackage{enumitem}
%\setlist{nosep}
%\setitemize{topsep=0.5em, itemsep=0.5em}
%\setenumerate{topsep=0.5em, itemsep=0.5em}
%\setitemize[1]{label=\bullet}
%\setitemize[2]{label=\circ}

% Adjust page margins
\usepackage{geometry} % Required for adjusting page dimensions and margins

\geometry{%
	paper=a4paper, % Change to letterpaper for US letter
	%total={209mm,157mm},
	top=3cm,% Top margin
	bottom=3cm,% Bottom margin
	left=2.5cm,% Left margin
	right=2.5cm, % Right margin
	headheight=30pt, % Header height
	footskip=1.4cm, % Space from the bottom margin to the baseline of the footer
	headsep=1.2cm, % Space from the top margin to the baseline of the header
	%showframe, % Uncomment to show how the type block is set on the page
}

% Don't indent paragraphs, leave some space between them
\usepackage{parskip}

% Floating elements
%\usepackage{wrapfig}

% Hide page number when page is empty
%\usepackage{emptypage}

\usepackage{caption}
\usepackage{subcaption}
\usepackage{multicol}
\usepackage[table,x11names]{xcolor}

% Hyper ref
\usepackage[
	pdfencoding=auto, psdextra, % Enable unicode and math in titles
	colorlinks=true,
	urlcolor=DodgerBlue4,
	linkcolor=DodgerBlue4,
	citecolor=DodgerBlue4
]{hyperref}
\PassOptionsToPackage{unicode}{hyperref}
\PassOptionsToPackage{naturalnames}{hyperref}
\pdfstringdefDisableCommands{\def\varepsilon{\textepsilon}}
\usepackage{bookmark}% faster updated bookmarks
%\usepackage{csquotes}

% Bibliography
% \usepackage[nottoc,numbib]{tocbibind} % add toc in index
%\usepackage[
%	autocite     = superscript,
%	backend      = biber,
%	sortcites    = true,
%	style        = numeric,
%	defernumbers = true
%]{biblatex}

% }}}

% {{{ Sectioning

%\usepackage{titlesec} % Required for modifying sections
\usepackage{tikz}

%------------------------------------------------
% Part

%\titleformat%
%{\part} % command
%[display] % shape
%{\bfseries\huge\scshape} % format
%{
%	\center%
%	Parte\ \thepart
%} % label
%{0.5ex} % sep
%{
%	\center
%} % before-code
%[
%	\vspace{-0.5ex}%
%] % after-code
%
%%------------------------------------------------
%% Chapter
%
%%\addto\captionsspanish{\renewcommand{\chaptername}{\assignmentChapterName}}
%%Options: Sonny, Lenny, Glenn, Conny, Rejne, Bjarne, Bjornstrup
%\titleformat%
%{\chapter} % command
%[hang] % shape
%{\LARGE\scshape} % format
%{\thechapter\ \ } % label
%{0.5ex} % sep
%{} % before-code
%[
%	\vspace{-0.5ex}%
%	\color{DodgerBlue3}\rule{\textwidth}{0.3pt}
%] % after-code
%
%\titlespacing{\chapter}{0pt}{-0.5cm}{1.5em}
%
%%------------------------------------------------
%% Section
%
%\titleformat%
%{\section} % command
%[hang] % shape
%{\bfseries\LARGE} % format
%{\thesection\ \ } % label
%{0.5ex} % sep
%{} % before-code
%[] % after-code
%
%
%%------------------------------------------------
%% SubSection
%
%\titleformat%
%{\subsection} % command
%[hang] % shape
%{\bfseries\Large} % format
%{\thesubsection\ \ } % label
%{0.5ex} % sep
%{} % before-code
%[] % after-code
%
%%------------------------------------------------
%% SubSubSection
%
%\titleformat%
%{\subsubsection}
%[hang]
%{\bfseries\large}
%{\thesubsubsection\ \ }
%{0.5ex}
%{}
%[]
%
%% }}}

% {{{ Math stuff

% Math stuff
\usepackage{amsmath}
\usepackage{amsfonts}
\usepackage{mathtools}
\usepackage{amsthm}
\usepackage{amssymb}
\usepackage{physics}

% Fancy script capitals
\usepackage{mathrsfs}
%\usepackage{cancel}

% Math code
%\usepackage{sagetex}

% Bold math
\usepackage{bm}

% Some shortcuts
\newcommand\N{\ensuremath{\mathbb{N} }}
\newcommand\M{\ensuremath{\mathfrak{M} }}
\newcommand\R{\ensuremath{\mathbb{R} }}
\newcommand\K{\ensuremath{\mathbb{K} }}
\newcommand\Z{\ensuremath{\mathbb{Z} }}
\renewcommand\O{\ensuremath{\emptyset}}
\newcommand\Q{\ensuremath{\mathbb{Q} }}
\newcommand\C{\ensuremath{\mathbb{C} }}
\newcommand\Dom{\ensuremath{\textrm{Dom}} }

% Add \contra symbol to denote contradiction
\usepackage{stmaryrd} % for \lightning
\newcommand\contra{\scalebox{1.5}{$\lightning$}}

% Theorem tools
\usepackage{thmtools}
%\usepackage{cleveref}

%\DeclareMathOperator{\sech}{sech}
%\DeclareMathOperator{\csch}{csch}
\DeclareMathOperator{\arccosh}{arcCosh}
\DeclareMathOperator{\arcsinh}{arcsinh}
\DeclareMathOperator{\arctanh}{arctanh}
\DeclareMathOperator{\arcsech}{arcsech}
\DeclareMathOperator{\arccsch}{arcCsch}
\DeclareMathOperator{\arccoth}{arcCoth}

%\newcommand\oast{\stackMath\mathbin{\stackinset{c}{0ex}{c}{0ex}{\ast}{\bigcirc}} }

% Theorems in fancy boxes
%\usepackage{todonotes}
%\usepackage[most]{tcolorbox}
%\tcbuselibrary{theorems}
%\tcbuselibrary{minted}
%\tcbset{listing engine=minted}

%\theoremstyle{definition}
%\newtheorem{definition}{Definition}[section]
%\newtcbtheorem[number within=section]{definition}{Definition}%
%{colback=green!5,colframe=green!35!black,fonttitle=\bfseries}{def}

\definecolor{codegreen}{rgb}{0,0.6,0}
\definecolor{codegray}{rgb}{0.5,0.5,0.5}
\definecolor{codepurple}{rgb}{0.58,0,0.82}
\definecolor{backcolour}{rgb}{0.95,0.95,0.92}

\lstdefinestyle{mystyle}{
	backgroundcolor=\color{backcolour},
	commentstyle=\color{codegreen},
	keywordstyle=\color{magenta},
	numberstyle=\tiny\color{codegray},
	stringstyle=\color{codepurple},
	basicstyle=\ttfamily\footnotesize,
	breakatwhitespace=false,
	breaklines=true,
	captionpos=b,
	keepspaces=true,
	numbers=left,
	numbersep=5pt,
	showspaces=false,
	showstringspaces=false,
	showtabs=false,
	tabsize=2
}

\lstset{style=mystyle}

%\newtcbinputlisting{\mathcodeinput}[3]{
%	listing engine=minted,
%	%minted style=colorful,
%	minted language={#2},
%	minted options={
%			fontsize=\small,
%			breaklines,
%			autogobble,
%			linenos,
%			numbersep=3mm,
%			mathescape=true
%		},
%	colback=DodgerBlue1!6,
%	colframe=DodgerBlue4,
%	boxrule=0.0mm,
%	left=5mm,
%	listing file={#1},
%	enhanced,
%	breakable,
%	overlay={
%			\begin{tcbclipinterior}
%				\fill[DodgerBlue2!20]
%				(frame.south west)rectangle([xshift=5mm]frame.north west);
%			\end{tcbclipinterior}
%		},
%	#3
%}

%\newtcblisting{mathcodebox}[2][]{
%	listing engine=minted,
%	%minted style=colorful,
%	minted language={#2},
%	minted options={
%			fontsize=\small,
%			breaklines,
%			autogobble,
%			linenos,
%			numbersep=3mm
%		},
%	colback=DodgerBlue1!6,
%	colframe=DodgerBlue4,
%	boxrule=0.0mm,
%	left=5mm,
%	enhanced,
%	breakable,
%	overlay={
%			\begin{tcbclipinterior}
%				\fill[DodgerBlue2!20]
%				(frame.south west)rectangle([xshift=5mm]frame.north west);
%			\end{tcbclipinterior}
%		},
%	#1
%}

% End example and intermezzo environments with a small diamond (just like proof
% environments end with a small square)
\usepackage{etoolbox}
%\AtEndEnvironment{eg}{\null\hfill$\qed$}%
\AtEndEnvironment{myproof}{\null\hfill$\square$}%

\newcommand{\exerciseSection}[1]{%
	\pagebreak
	\section*{#1}
	\addcontentsline{toc}{section}{\protect\numberline{}#1}%
}

%\usepackage{xifthen}

% Exercise
% Usage:
% \exercise{5}
% \subexercise{1}
% \subexercise{2}
% \subexercise{3}
% gives
% Problema 5
%   Problema 5.1
%   Problema 5.2
%   Problema 5.3
%\newcommand{\exercise}[1]{%
%  \problemDivider
%  \problem{#1}
%  \def\@exercise{#1}%
%}
%
%\newcommand{\subexercise}[1]{%
%  \problem{\@exercise.#1}
%}

% }}}

% {{{ Physics

\usepackage{physics}

% }}}

% {{{ Class stuff

% \lecture starts a new lecture (les in dutch)
%
% Usage:
% \lecture{1}{di 12 feb 2019 16:00}{Inleiding}
%
% This adds a section heading with the number / title of the lecture and a
% margin paragraph with the date.

% I use \dateparts here to hide the year (2019). This way, I can easily parse
% the date of each lecture unambiguously while still having a human-friendly
% short format printed to the pdf.

\def\testdateparts#1{\dateparts#1\relax}
\def\dateparts#1 #2 #3 #4 #5\relax{
	\marginpar{\small\textsf{\mbox{#1 #2 #3 #5}} }
}

\def\@lecture{}%
\newcommand{\lecture}[3]{
	\ifthenelse{\isempty{#3}}{%
		\def\@lecture{Class #1}%
	}{%
		\def\@lecture{Class #1: #3}%
	}%
	\subsection*{\@lecture}
	\marginpar{\small\textsf{\textbf{Class #1} \\ \mbox{#2}} }
}

\newcommand{\lectureMargin}[2]{
	\marginpar{\small\textsf{\textbf{Clase #1} \\ \mbox{#2}} }
}

\def\@studySession{}%
\newcommand{\studySession}[3]{
	\ifthenelse{\isempty{#3}}{%
		\def\@studySession{Study #1}%
	}{%
		\def\@studySession{Study #1: #3}%
	}%
	\subsection*{\@studySession}
	\marginpar{\small\textsf{\def\@studySession{Study #1}\\ \mbox{#2}} }
}

\newcommand{\studySessionMargin}[2]{
	\marginpar{\small\textsf{\textbf{Study #1} \\ \mbox{#2}} }
}

%\usepackage{bbding}
%\newcommand{\pendingCorrectionMargin}[2]{
%	\marginpar{\footnotesize\textsf{\textbf{\PencilLeftDown\, #1} \\ \mbox{#2}} }
%}

%\newenvironment{mathProblem}{
%  \vspace{0.5\baselineskip} % Whitespace before the question
%  \paragraph{} % Blank section title (e.g. just Question 2)
%  \lfoot{\small\itshape\assignmentQuestionName~\thesection~continuada en la siguiente página\ldots} % Set the left footer to state the question continues on the next page, this is reset to nothing if it doesn't (below)
%}{
%  \lfoot{} % Reset the left footer to nothing if the current question does not continue on the next page
%}

% }}}

% {{{ Header and footer

% LE: left even
% RO: right odd
% CE, CO: center even, center odd
% My name for when I print my lecture notes to use for an open book exam.
% \fancyhead[LE,RO]{Gilles Castel}

\makeatletter
\let\runauthor\@author
\let\runtitle\@title
\makeatother

%\newcommand{\uniIcon}{../../../utad-logo.png}
\newcommand{\uniIcon}{assets/img/utad-icon.png}
%\newcommand{\uniIcon}{../../../utad-icon.png}

\usepackage{fancyhdr} % Required for customizing headers and footers
\usepackage{refcount}

\pagestyle{fancy} % Enable custom headers and footers

\renewcommand\headrulewidth{0.5pt} % Thickness of the header rule

\fancyhead{} % Clear default header
\fancyfoot{} % Clear default footer

\fancyhead[L]{\small\runtitle}
\fancyhead[C]{\includegraphics[height=0.85em, trim= 0 0 0 -1em]{\uniIcon}}
\fancyhead[R]{\small\runauthor}

\fancyfoot[L]{\small{\leftmark}}
%\fancyfoot[R]{\small \textcolor{gray}{PÁGINA} \thepage\ | \getrefbykeydefault{LastPage}{page}{0}}% Número de página

\renewcommand{\headrulewidth}{0.5pt}% 1pt header rule
\renewcommand{\headrule}{\hbox to\headwidth{%
		\color{DodgerBlue4}\leaders\hrule height\headrulewidth\hfill}}
\renewcommand{\footrulewidth}{0pt}% No footer rule

\fancypagestyle{plain}{
	\fancyhead[L]{\small\runtitle}
	\fancyhead[C]{\includegraphics[height=0.85em, trim= 0 0 0 -1em]{\uniIcon}}
	\fancyhead[R]{\small\runauthor}

	\fancyfoot[L]{\small{\leftmark}}
	\fancyfoot[R]{
		\small \textcolor{gray}{PAGE} \thepage\ | \getrefbykeydefault{LastPage}{page}{0}}% Número de página

	\renewcommand{\headrulewidth}{0.5pt}% 1pt header rule
	\renewcommand{\headrule}{\hbox to\headwidth{%
			\color{DodgerBlue4}\leaders\hrule height\headrulewidth\hfill}}
	\renewcommand{\footrulewidth}{0pt}% No footer rule
}

\setlength{\headheight}{24pt} % Suppress warnings

\makeatother

% }}}

% {{{ Notes


% Corrección is correction in Spanish
% Usage:
% \begin{verbetering}
%     Lorem ipsum dolor sit amet, consetetur sadipscing elitr, sed diam nonumy eirmod
%     tempor invidunt ut labore et dolore magna aliquyam erat, sed diam voluptua. At
%     vero eos et accusam et justo duo dolores et ea rebum. Stet clita kasd gubergren,
%     no sea takimata sanctus est Lorem ipsum dolor sit amet.
% \end{verbetering}
\newenvironment{correction}{\begin{tcolorbox}[
			arc=0mm,
			colback=white,
			colframe=green!60!black,
			title=Opmerking,
			fonttitle=\sffamily,
			breakable
		]}{\end{tcolorbox}}

% Nota is note in Dutch. Same as 'verbetering' but color of box is different
%\newenvironment{note}[1]{\begin{tcolorbox}[
%      arc=0mm,
%      colback=white,
%      colframe=white!60!black,
%      title=#1,
%      fonttitle=\sffamily,
%      breakable
%    ]}{\end{tcolorbox}}

\makeatletter
\AtBeginDocument{%
	\newcommand\My@Macro[1]{#1}%
	\newcommand\My@Thirdoffive[5]{\My@Macro{#3}}%
	\renewcommand*\@namerefstar[1]{%
		\HyRef@StarSetRef{#1}\My@Thirdoffive
	}%
	\renewcommand*\T@nameref[1]{%
		\begingroup
		\let\label\@gobble
		\NR@setref{#1}\My@Thirdoffive{#1}%
		\endgroup
	}%
	\DeclareRobustCommand\fucnameref{%
		\@ifstar\fucnameref@star\fucnameref@nostar
	}%
	\newcommand\callemakefirstuc[1]{%
		\MakeLowercase{\emakefirstuc{#1}}%
	}%
	\newcommand\fucnameref@star[1]{%
		\begingroup
		\let\My@Macro=\callemakefirstuc
		\nameref*{#1}%
		\endgroup
	}%
	\newcommand\fucnameref@nostar[1]{%
		\begingroup
		\let\My@Macro=\callemakefirstuc
		\nameref{#1}%
		\endgroup
	}%
	\DeclareRobustCommand\ucnameref{%
		\@ifstar\ucnameref@star\ucnameref@nostar
	}%
	\newcommand\ucnameref@star[1]{%
		\begingroup
		\MFUhyphentrue
		\let\My@Macro=\ecapitalisefmtwords
		\nameref*{#1}%
		\endgroup
	}%
	\newcommand\ucnameref@nostar[1]{%
		\begingroup
		\MFUhyphentrue
		\let\My@Macro=\ecapitalisefmtwords
		\nameref{#1}%
		\endgroup
	}%
	\DeclareRobustCommand\lcnameref{%
		\@ifstar\lcnameref@star\lcnameref@nostar
	}%
	\newcommand\lcnameref@star[1]{%
		\begingroup
		\let\My@Macro=\MakeLowercase
		\nameref*{#1}%
		\endgroup
	}%
	\newcommand\lcnameref@nostar[1]{%
		\begingroup
		\let\My@Macro=\MakeLowercase
		\nameref{#1}%
		\endgroup
	}%
}%
\makeatother

% }}}

% {{{ Figure support

% Figure support as explained in my blog post.
%\usepackage{import}
%\usepackage{xifthen}
%\usepackage{pdfpages}
%\usepackage{transparent}
\newcommand{\incfig}[1]{%
	\def\svgwidth{\columnwidth}
	%\import{./figures/}{#1.pdf_tex}
}

%\usepackage{svg}

% }}}

% {{{ Diagram support
%\usepackage{pgfplots}
%\pgfplotsset{colormap={utad}{color=(DodgerBlue1) color=(black) color=(Firebrick1)}}
%\pgfplotsset{
%	compat=newest,
%	colormap name=utad
%}
%\usepgfplotslibrary{polar}
%\usetikzlibrary{
%	3d,
%	angles,
%	arrows,
%	babel,
%	calc,
%	decorations.markings,
%	patterns,
%	pgfplots.fillbetween,
%	quotes
%}
% }}}
